\section{Validation}

% --- Comparaison avec le livre ---
Le script refait d'abord les premiers pas de l'exemple du livre, avec $\tau$ arrondi à
\SI{0.55}{ms} comme dans la solution.

\begin{table}[H]
\centering
\caption{Charge R-L : valeurs du livre et du script ($\tau = \SI{0.55}{ms}$).}
\label{tab:livre}
\begin{tabular}{llrrrr}
\toprule
$t$ (ms) & source & $i_1$ (A) & $v_1$ (V) & $i_2$ (A) & $v_2$ (V) \\
\midrule
0,1 & livre  & \livreiunA & \livrevunA & 0 & 0 \\
    & script & \iunA & \vunA & \ideuxA & \vdeuxA \\
0,2 & livre  & \livreiunB & \livrevunB & 0 & 0 \\
    & script & \iunB & \vunB & \ideuxB & \vdeuxB \\
0,7 & livre  & \livreiunC & \livrevunC & \livreideuxC & \livrevdeuxC \\
    & script & \iunC & \vunC & \ideuxC & \vdeuxC \\
\bottomrule
\end{tabular}
\end{table}

Les valeurs coïncident à l'arrondi près. Le seul écart visible porte sur $v_2$ à
\SI{0.7}{ms} : \vdeuxC~V contre \livrevdeuxC~V. Le livre arrondit $Z_c$ à \SI{274}{\ohm} et
garde des valeurs intermédiaires arrondies (la source d'historique $I_m$ est écrite 1368~A),
alors que le script garde $Z_c = \Zc~\unit{\ohm}$ et toutes les décimales. L'onde arrive en $m$
au même pas que dans le livre, \SI{0.7}{ms}. La suite du calcul utilise le $\tau$ exact,
$\tauL~\unit{\micro s}$.

% --- Régime permanent par phaseurs ---
Le régime permanent de Q1 se calcule à la main en écrivant la ligne comme un quadripôle ABCD
sans pertes ($A = D = \cos\beta d$, $B = jZ_c\sin\beta d$, $C = j\sin\beta d/Z_c$, avec
$\beta d = \omega\tau$) fermé sur $Z_L = R + jL\omega$. Le tableau~\ref{tab:phaseurs} compare
ces amplitudes aux crêtes simulées entre 4,9 et 5~s.

\begin{table}[H]
\centering
\caption{Q1 : amplitudes en régime permanent, simulation et phaseurs.}
\label{tab:phaseurs}
\begin{tabular}{lrr}
\toprule
grandeur & simulation & phaseurs \\
\midrule
$|I_1|$ (A)  & \IunSim & \IunPh \\
$|I_2|$ (A)  & \IdeuxSim & \IdeuxPh \\
$|V_2|$ (kV) & \VdeuxSim & \VdeuxPh \\
\bottomrule
\end{tabular}
\end{table}

L'accord est au dixième près. Ce contrôle vérifie aussi le signe de $I_k$ : l'onde réfléchie ne
revient en $k$ qu'après $2\tau$, elle n'apparaît pas dans le tableau~\ref{tab:livre}, mais une
erreur de signe fausserait complètement le régime permanent.
